\documentclass[12pt]{article}
\usepackage{cocina1}

\begin{document}
\begin{ficha}{Servir}{Cortar reparte, grapar junta}{1}

\begin{encargo}
La pizza viene cortada de otra manera: prepara \textbf{justo} lo que pide el cliente.
\end{encargo}

\bloque{1}{Mira cómo se hace}{Con pizza}

\vspace{2mm}
\begin{minipage}[c]{.5\linewidth}
\paso{1}\textbf{Mira lo que hay.}\\[1mm]
\hspace*{6mm}Tres trozos de cuatro:\ \pizza[escala=.7]{4}{0,1,2}
\end{minipage}%
\begin{minipage}[c]{.5\linewidth}
\paso{2}\textbf{Corta cada trozo en 2.}\\[1mm]
\hspace*{6mm}Seis trozos de ocho:\ \pizza[escala=.7]{8}{0,1,2,3,4,5}
\end{minipage}

\vspace{2mm}
\begin{minipage}[c]{.5\linewidth}
\paso{3}\textbf{Escríbelo.}\quad \cortado{3}{4}{2}
\end{minipage}%
\begin{minipage}[c]{.5\linewidth}
{\small Compruebo: en la regla caen en el mismo sitio \bien}\\[1mm]
\regla[6.5cm]{1}{8}{3/4/tomate}
\end{minipage}

\vspace{1mm}
\begin{nick}
Cortar cada trozo en 2 \textbf{multiplica por 2} el número de arriba y el de
abajo. Hay más trozos, pero la pizza es \textbf{la misma}.
\end{nick}

\bloque{2}{Ahora tú}{Con pizza}\quad{\small dibuja los cortes por las rayas de puntos}

\vspace{1mm}
\begin{minipage}[t]{.5\linewidth}\vspace{0pt}
\textbf{a)} Viene en mitades; pide cuartos.\\[.5mm]
\pizza[escala=.7,guias=4]{2}{0}\quad\cortadohueco{1}{2}{2}
\end{minipage}%
\begin{minipage}[t]{.5\linewidth}\vspace{0pt}
\textbf{b)} Viene en tercios; pide sextos.\\[.5mm]
\pizza[escala=.7,guias=6]{3}{0,1}\quad\cortadohueco{2}{3}{2}
\end{minipage}

\bloque{3}{Al revés: grapar}{Con pizza}\quad{\small grapar de 2 en 2 \textbf{divide} entre 2}

\vspace{1mm}
\begin{minipage}[t]{.5\linewidth}\vspace{0pt}
\textbf{a)} Viene en cuartos; pide mitades.\\[.5mm]
\pizza[escala=.7,grapas={0}]{4}{0,1}\quad\grapadohueco{2}{4}{2}
\end{minipage}%
\begin{minipage}[t]{.5\linewidth}\vspace{0pt}
\textbf{b)} Viene en sextos; pide tercios.\\[.5mm]
\pizza[escala=.7,grapas={0,2}]{6}{0,1,2,3}\quad\grapadohueco{4}{6}{2}
\end{minipage}

\alto{cortar y grapar no suman nada. $\frac{1}{2}$ no se convierte en
$\frac{2}{3}$: eso ya es otra pizza.}

\reto{¿Cuántos octavos son media pizza? ¿Y cuántos doceavos? Explica
cómo lo sabes.}

\comprueba{Cada código abre ese pedido.}{%
\qr{1}{j=servir\&f=6/8\&c=4}\hfill\qr{2a}{j=servir\&f=2/4\&c=2}\hfill
\qr{2b}{j=servir\&f=4/6\&c=3}\hfill\qr{3a}{j=servir\&f=1/2\&c=4\&p=1}\hfill
\qr{3b}{j=servir\&f=2/3\&c=6}}

\end{ficha}

% ─────────────────────────────── REVERSO ──────────────────────────────────
\begin{dorso}{Servir}{Más pedidos}{1}

\bloque{1}{Sin pizza}{Sobre el papel}\quad{\small multiplica o divide arriba y abajo por el mismo número}

\vspace{2mm}
\textbf{a)}\ \cortadohueco{1}{3}{2}\hfill\textbf{b)}\ \cortadohueco{3}{5}{2}\hfill\textbf{c)}\ \cortadohueco{2}{4}{2}

\vspace{3mm}
\textbf{d)}\ \grapadohueco{6}{8}{2}\hfill\textbf{e)}\ \grapadohueco{5}{10}{5}\hfill\textbf{f)}\ \grapadohueco{3}{9}{3}

\begin{minipage}[t]{.6\linewidth}\vspace{0pt}
\bloque{2}{La regla del tique}{Con la regla}

\vspace{2mm}
Marca $\dfrac{1}{2}$, $\dfrac{2}{4}$ y $\dfrac{4}{8}$.\\[1mm]
\reglavacia[9.4cm]{1}{8}\\[1mm]
¿Qué pasa? \hueco[6cm]

\bloque{3}{Problemas}{Sobre el papel}

\vspace{1.5mm}
\textbf{a)} Nick corta una pizza en 8 trozos y Ana se come 4.
¿Qué fracción se ha comido? \hueco[1.4cm]\ ¿Es media pizza? \hueco[1.4cm]

\vspace{3mm}
\textbf{b)} Un cliente quiere $\frac{2}{3}$ de pizza y viene en sextos.
¿Cuántos trozos le das? \hueco[1.4cm]
\end{minipage}\hfill
\begin{minipage}[t]{.37\linewidth}\vspace{4mm}
\tablamult
\end{minipage}

\vspace{3mm}
\begin{semaforo}
\sfila{Corto los trozos y escribo la fracción}
\sfila{Grapo los trozos y escribo la fracción}
\sfila{Sé si dos fracciones son la misma pizza}
\end{semaforo}

\comprueba{Cada código abre ese pedido.}{%
\qr{1a}{j=servir\&f=2/6\&c=3}\hfill\qr{1b}{j=servir\&f=6/10\&c=5}\hfill
\qr{1d}{j=servir\&f=3/4\&c=8}\hfill\qr{1e}{j=servir\&f=1/2\&c=10\&p=1}\hfill
\qr{3a}{j=servir\&f=1/2\&c=8\&p=1}\hfill\qr{3b}{j=servir\&f=2/3\&c=6}}
\end{dorso}
\end{document}
